\documentclass[../../main.tex]{subfiles}% 注意这里的文件路径不能用 ./main.tex ,否则用latexmk编译子文件会报错
\graphicspath{{\subfix{./image/}}{\subfix{../../image/}}} % 指定图片引用路径,后续可以直接使用图片文件名
% 如果图片存放在上级目录,则使用 ../../image/ 作为备用路径

% 例如:
% \begin{figure}[H]
% \centering
% \includegraphics[scale=0.3]{图.png}
% \caption{}
% \label{figure:图}
% \end{figure}
% 注意:上述\label{}一定要放在\caption{}之后,否则引用图片序号会只会显示??.

\begin{document}

\section{主理想整环上的矩阵}

交换幺环上有限秩自由模的模同态可由环上矩阵刻画，因此可通过矩阵研究主理想整环上有限生成模的不变因子、初等因子、秩等不变量。

主理想整环上有限生成模为有限秩自由模的同态像，即自由模的商模。设$D$为主理想整环，$\mathscr{M}$为$D$上秩$m$的自由模，$\mathscr{N}\subseteq\mathscr{M}$为秩$n$的子模。取$\mathscr{M}$的基$e_1,e_2,\cdots,e_m$，$\mathscr{N}$的基$f_1,f_2,\cdots,f_n$，则存在$m\times n$矩阵$A\in M_{m\times n}(D)$满足
$$(f_1,f_2,\cdots,f_n)=(e_1,e_2,\cdots,e_m)A.$$
若更换两组基：
$$(e_1,e_2,\cdots,e_m)=(e_1',e_2',\cdots,e_m')P,\quad (f_1,f_2,\cdots,f_n)=(f_1',f_2',\cdots,f_n')Q,$$
其中$P\in GL_m(D)$，$Q\in GL_n(D)$，则
$$(f_1',f_2',\cdots,f_n')=(e_1',e_2',\cdots,e_m')PAQ^{-1}.$$

\begin{definition}
设$D$为主理想整环，$A,B\in M_{m\times n}(D)$。若存在$P\in GL_m(D)$，$Q\in GL_n(D)$使得$B=PAQ$，则称$A$与$B$等价。
\end{definition}

\begin{definition}
记$D^*=D\setminus\{0\}$，定义函数$l\colon D^*\to\mathbb{N}\cup\{0\}$：
\begin{align*}
l(a)=
\begin{cases}
0,&a\text{可逆},\\
s,&a=p_1p_2\cdots p_s,\ p_i\text{为素元}.
\end{cases}
\end{align*}
\end{definition}

\begin{proposition}
函数$l$满足如下性质：
\begin{enumerate}[(1)]
\item 若$a,b$相伴，则$l(a)=l(b)$；
\item 若$a\mid b$且$b\nmid a$，则$l(a)<l(b)$；
\item 若$a\mid b$，则$l(a)=l(b)$当且仅当$a,b$相伴。
\end{enumerate}
\end{proposition}

\begin{definition}
设$E_{ij}$为$(i,j)$元为$1$、其余元为$0$的方阵，即$\mathrm{ent}_{kl}(E_{ij})=\delta_{ki}\delta_{lj}$。定义三类方阵：
\begin{align*}
P(i,j)&=E_{ij}+E_{ji}+\sum_{k\neq i,j}E_{kk},\\
P(i(c))&=cE_{ii}+\sum_{k\neq i}E_{kk},\quad c\in D^\times,\\
P(j,i(k))&=I+kE_{ji},\quad I\text{为单位矩阵}.
\end{align*}
\end{definition}

\begin{remark}
上述三类矩阵均可逆，左（右）乘矩阵对应矩阵的初等行（列）变换，故矩阵经初等变换后与原矩阵等价。
\end{remark}

\begin{theorem}
设$D$为主理想整环，$A\in M_{m\times n}(D)$，则$A$等价于形如
\begin{align*}
\begin{pmatrix}
d_1&&&\\
&d_2&&\\
&&\ddots&\\
&&&d_r\\
&&&&\boldsymbol{0}
\end{pmatrix}
\end{align*}
的矩阵，满足$d_i\mid d_{i+1}\ (1\leqslant i\leqslant r-1)$。该矩阵称为$A$的标准形，$d_1,d_2,\cdots,d_r$称为$A$的不变因子。
\end{theorem}

\begin{proof}

若$A=\boldsymbol{0}$，结论显然成立。下设$A\neq\boldsymbol{0}$。

第一步，选取$a_{kl}\neq0$满足$l(a_{kl})\leqslant l(a_{ij})$对所有$a_{ij}\neq0$成立。通过行、列互换，可使等价矩阵的$(1,1)$元为$a_{kl}$，仍记为$A$，则$a_{11}\neq0$且$l(a_{11})\leqslant l(a_{ij})$对所有$a_{ij}\neq0$成立。

第二步，构造可逆矩阵$Q$，使得$AQ$满足$a_{11}\mid a_{1k},\ a_{11}\mid a_{l1}\ (2\leqslant k\leqslant n,\ 2\leqslant l\leqslant m)$。
若$a_{11}\nmid a_{1k}$，设$d=(a_{11},a_{1k})$，则$a_{11}=a_1'd,\ a_{1k}=a_k'd$，存在$u,v\in D$使得$ua_{11}+va_{1k}=1$。令
\begin{align*}
Q=\left(\sum_{j\neq1,k}E_{jj}\right)+uE_{11}+a_1'E_{kk}+vE_{k1}-a_k'E_{1k},
\end{align*}
则
\begin{align*}
Q^{-1}=\left(\sum_{j\neq1,k}E_{jj}\right)+a_1'E_{11}+uE_{kk}-vE_{k1}+a_k'E_{1k}.
\end{align*}
此时$\mathrm{ent}_{11}(AQ)=d$，$\mathrm{ent}_{1k}(AQ)=0$，且$l(d)<l(a_{11})$。对行作同理处理，有限步后可得$a_{11}\mid a_{1k},\ a_{11}\mid a_{l1}$。

若$a_{11}\nmid a_{ij}$，设$a_{i1}=a_{11}q$，将第$i$行加第$1$行的$-q$倍，得等价矩阵$A_1$，满足$\mathrm{ent}_{i1}(A_1)=0$。再将第$i$行加到第$1$行得$A_2$，此时$a_{11}\nmid\mathrm{ent}_{1j}(A_2)$。若$l(a_{11})>l(\mathrm{ent}_{1j}(A_2))$，回到第一步；若$l(a_{11})<l(\mathrm{ent}_{1j}(A_2))$，回到第二步。因$l(a)\geqslant0$，该过程有限终止。

故有限步后，$a_{11}\mid a_{ij}$对所有$i,j$成立。将第$i$行减去第$1$行的$\frac{a_{i1}}{a_{11}}$倍，第$j$列减去第$1$列的$\frac{a_{1j}}{a_{11}}$倍，可得等价矩阵
\begin{align*}
\begin{pmatrix}
a_{11}&\\
&B
\end{pmatrix},\quad B\in M_{(m-1)\times(n-1)}(D),\ a_{11}\mid\mathrm{ent}_{ij}(B).
\end{align*}
对$B$归纳使用上述过程，可得$A$等价于标准形。

\end{proof}

\begin{corollary}
设$D$为主理想整环，$\mathscr{M}$为秩$m$的自由$D$-模，$\mathscr{N}\subseteq\mathscr{M}$为子模，则存在$\mathscr{M}$的基$e_1,e_2,\cdots,e_m$，以及$D$中非零元$d_1,d_2,\cdots,d_r$满足：
\begin{enumerate}[(1)]
\item $d_i\mid d_{i+1}\ (1\leqslant i\leqslant r-1)$；
\item $d_1e_1,d_2e_2,\cdots,d_re_r$为$\mathscr{N}$的一组基，称$d_1,\cdots,d_r$为$\mathscr{N}$的不变因子。
\end{enumerate}
\end{corollary}

\begin{proof}

任取$\mathscr{M}$的基$e_1',\cdots,e_m'$，$\mathscr{N}$的生成元$f_1',\cdots,f_n'$，则存在$A_1\in M_{m\times n}(D)$使得$f_j'=(e_1',\cdots,e_m')\mathrm{col}_jA_1$。

由前述定理，存在$P\in GL_m(D)$，$Q\in GL_n(D)$使得$A=PA_1Q$为标准形。令
$$(e_1,\cdots,e_m)=(e_1',\cdots,e_m')P^{-1},\quad (f_1,\cdots,f_n)=(f_1',\cdots,f_n')Q,$$
则$e_1,\cdots,e_m$为$\mathscr{M}$的基，$f_1,\cdots,f_n$为$\mathscr{N}$的生成元，且
$$(f_1,\cdots,f_n)=(e_1,\cdots,e_m)A.$$
故$f_i=d_ie_i\ (1\leqslant i\leqslant r)$，$f_i=\boldsymbol{0}\ (i>r)$。由$e_1,\cdots,e_m$线性无关，得$d_1e_1,\cdots,d_re_r$线性无关，即为$\mathscr{N}$的基。

\end{proof}

\begin{remark}
由上述定理与推论，有模同构
$$\mathscr{M}/\mathscr{N}\cong\bigoplus\limits_{i=1}^m(De_i/Dd_ie_i),$$
其中$i>r$时$d_i=0$。该结论可由自由模基本定理得到。
\end{remark}

\begin{remark}
当$D$为欧氏环时，可用赋值函数$\delta(a)$替代$l(a)$，辗转相除法中余数可替代最大公因子简化计算。
\end{remark}

\begin{definition}
设$D$为主理想整环，$A\in M_{m\times n}(D)$，$k\geqslant1$。$A$的所有$k$阶子式的最大公因子称为$A$的$k$级行列式因子，记为$D_k(A)$，规定$D_0(A)=1$。记$A\begin{pmatrix}i_1i_2\cdots i_k\\j_1j_2\cdots j_k\end{pmatrix}$为划去第$i_1,\cdots,i_k$行、第$j_1,\cdots,j_k$列的$k$阶子式。

若存在$r\in\mathbb{N}$使得$D_r(A)\neq0$且$D_{r+1}(A)=0$，则称$r$为$A$的秩，记为$\mathrm{rank}A=r$。当$k>r$时，$D_k(A)=0$。
\end{definition}

\begin{remark}
由定义可得$D_k(A)\mid D_{k+1}(A)\ (k=0,1,\cdots)$。
\end{remark}

\begin{theorem}
设$D$为主理想整环，$A,B\in M_{m\times n}(D)$。若$A$与$B$等价，则$D_k(A)$与$D_k(B)$相伴。
\end{theorem}

\begin{proof}

设$B=PAQ$，$P\in GL_m(D)$，$Q\in GL_n(D)$。令$C=PA$，先证$A$与$C$行列式因子相伴。

对任意$k$阶子式，$\det C\begin{pmatrix}i_1\cdots i_k\\j_1\cdots j_k\end{pmatrix}=\det(P_kA_k)$，其中$P_k$为$P$的$i_1,\cdots,i_k$行构成的矩阵，$A_k$为$A$的$j_1,\cdots,j_k$列构成的矩阵。将行列式按行展开，可得$D_k(A)\mid\det C\begin{pmatrix}i_1\cdots i_k\\j_1\cdots j_k\end{pmatrix}$，故$D_k(A)\mid D_k(C)$。

又$A=P^{-1}C$，同理$D_k(C)\mid D_k(A)$，故$D_k(A)\sim D_k(C)$。

对$B=CQ$转置得$B'=Q'C'$，同理$D_k(B')\sim D_k(C')$，即$D_k(B)\sim D_k(C)$。故$D_k(A)\sim D_k(B)$。

\end{proof}

\begin{theorem}
设$D$为主理想整环，$A\in M_{m\times n}(D)$，$d_1,\cdots,d_r$为$A$的不变因子，$D_k(A)$为$k$级行列式因子，则
\begin{enumerate}[(1)]
\item $d_k\sim\frac{D_k(A)}{D_{k-1}(A)}\ (1\leqslant k\leqslant r)$；
\item $d_1,\cdots,d_r$在相伴意义下唯一；
\item $B\in M_{m\times n}(D)$与$A$等价当且仅当$A,B$有相同的不变因子（或等价地，相同的行列式因子）。
\end{enumerate}
\end{theorem}

\begin{proof}

(1) 设$PAQ$为$A$的标准形，由前述定理，$D_k(A)\sim D_k(PAQ)$。结合$d_i\mid d_{i+1}$，得
\begin{align*}
D_k(A)=d_1d_2\cdots d_k,\quad 1\leqslant k\leqslant r,\\
D_k(A)=0,\quad k>r.
\end{align*}
故$d_k\sim\frac{D_k(A)}{D_{k-1}(A)}$。

(2) 若$P_1AQ_1$为另一标准形，不变因子为$d_1',\cdots,d_r'$，则$D_k(P_1AQ_1)\sim D_k(PAQ)$，故$d_k'\sim d_k$，即不变因子相伴唯一。

(3) 若$A,B$等价，则行列式因子相伴，故不变因子相同。若不变因子相同，则$A,B$等价于同一标准形，故等价。

\end{proof}

\begin{theorem}
设$V$为域$F$上$n$维线性空间，线性变换$\mathscr{A}$在基$u_1,\cdots,u_n$下的矩阵为$A$。将$V$视为$F[\lambda]$-模，$\eta\colon F[\lambda]^{(n)}\to V$为模同态满足$\eta(e_i)=u_i$，$N=\ker\eta$，则$N$有基
$$f_i=\lambda e_i-(e_1,\cdots,e_n)\mathrm{col}_iA,\quad 1\leqslant i\leqslant n.$$
\end{theorem}

\begin{proof}

首先，
\begin{align*}
\eta(f_i)&=\lambda\eta(e_i)-(\eta(e_1),\cdots,\eta(e_n))\mathrm{col}_iA\\
&=\mathscr{A}(u_i)-(u_1,\cdots,u_n)\mathrm{col}_iA=\boldsymbol{0},
\end{align*}
故$f_i\in N$。

记$A_1e_i=(e_1,\cdots,e_n)\mathrm{col}_iA$，则$\lambda e_i=f_i+A_1e_i$。对任意$\sum_{i=1}^ng_i(\lambda)e_i\in F[\lambda]^{(n)}$，存在$h_i(\lambda)\in F[\lambda]$，$b_i\in F$使得
$$\sum_{i=1}^ng_i(\lambda)e_i=\sum_{i=1}^nh_i(\lambda)f_i+\sum_{i=1}^nb_ie_i.$$
若$\sum_{i=1}^ng_i(\lambda)e_i\in N$，则$\sum_{i=1}^nb_iu_i=\boldsymbol{0}$，故$b_i=0$，即$\{f_i\}$生成$N$。

下证$\{f_i\}$线性无关。若$\sum_{i=1}^nh_i(\lambda)f_i=\boldsymbol{0}$，则
$$\sum_{i=1}^nh_i(\lambda)\lambda e_i=\sum_{i=1}^nh_i(\lambda)(e_1,\cdots,e_n)\mathrm{col}_iA.$$
由$e_1,\cdots,e_n$为$F[\lambda]^{(n)}$的基，得
$$\lambda h_i(\lambda)=(\mathrm{row}_iA)\begin{pmatrix}\lambda h_1(\lambda)\\\vdots\\\lambda h_n(\lambda)\end{pmatrix}.$$
取次数最高的$h_r(\lambda)$，比较次数可知上式矛盾，故$h_i(\lambda)=0$，即$\{f_i\}$线性无关，为$N$的基。

\end{proof}

\begin{remark}
求$F[\lambda]$-模$V$的不变因子等价于求$\lambda I_n-A$的不变因子，由此可建立线性变换理论与主理想整环上有限生成模理论的联系。
\end{remark}



\end{document}